%%%%%%%%%%%%%%%%%%%%%%%%%%%%%%%%%%%%%%%%%%%%%%%%%%%%%%
\documentclass[12pt,a4paper]{report}
% o article, book, ...

% TODO inserire vari packages (todonotes, ecc.)

%%%%%%%%%%%%%%%%%%%%%%%%%%%%%%%%%%%%%%%%%%%%%%%%%%%%%%
% packages...
\usepackage[utf8]{inputenc}
\usepackage[english,italian]{babel}
\usepackage[hyphens]{url}

% Per generare il file PDF aderente alle specifiche PDF/A-1b. Verificarne poi la validità.
%\usepackage[a-1b]{pdfx}

\usepackage{hyperref}

\usepackage{graphicx}

\usepackage{lipsum}% Per inserire testo a caso in attesa di realizzare i capitoli

\usepackage{listings}
\lstset{
	% 	language=bash
	frame=single,
	breaklines=true,
	postbreak=\raisebox{0ex}[0ex][0ex]{\ensuremath{\color{red}\hookrightarrow\space}},
	basicstyle=\ttfamily\footnotesize
}

\usepackage[style=alphabetic,
%citestyle=authoryear
]{biblatex}
\addbibresource{Biblio.bib}

\usepackage{epigraph} % per le frasi inizio capitolo

\usepackage{fancyhdr}

\usepackage[colorinlistoftodos]{todonotes}
%\usepackage[disable]{todonotes}

\usepackage{refcheck}


%%%%%%%%%%%%%%%%%%%%%%%%%%%%%%%%%%%%%%%%%%%%%%%%%%%%%
\begin{document}

% Frontespizio
\begin{titlepage}
\begin{center}
\includegraphics[width=\textwidth]{Logo.jpg}\\
{\large{\bf Corso di Laurea in Informatica}}
\end{center}
\vspace{12mm}
\begin{center}
{\huge{\bf Black-box per flotte a noleggio: monitoraggio/tracking veicoli, segnalazione allarmi e profilazione guida}}\\
\end{center}
\vspace{12mm}
\begin{flushleft}
{\large{\bf Relatore: Prof. Andrea Trentini}}
{\large{}}\\
\vspace{4mm}
{\large{\bf Correlatore: Dott. Alexjan Carraturo}}
{\large{}}\\
\end{flushleft}
\vspace{12mm}
\begin{flushright}
{\large{\bf Tesi di Laurea di:}}
{\large{Alessandro Belfiore}}\\
{\large{\bf Matr. 975826}}\\
\end{flushright}
\vspace{4mm}
\begin{center}
{\large{\bf Anno Accademico 2023/2024}}
\end{center}
\end{titlepage}


\tableofcontents


% o sections (dipende dal documentclass)
\chapter{Introduzione}
\section{Cos'è una scatola nera?}
\subsection{Utilità}
\subsection{Importanza}
\section{Obiettivi del progetto}

TCP/IP over Avian Carriers\cite{waitzman1990standard}

\chapter{Stato dell'arte}
\section{Tecnologie attuali di monitoraggio e tracking per veicoli}
\section{Gestione delle emergenze e sistemi di allarmi in uso}
\section{Soluzioni esistenti per la profilazione guidatori}

\chapter{Funzionalità del sistema}
\section{Monitoraggio}
\section{Tracking}
\section{Segnalazione di allarmi}
\subsection{Incidente}
\subsection{Assistenza}
\subsection{Rilevamento allarmi}
\subsubsection{Critici}
\subsubsection{Non critici}
\section{Profilazione del comportamento del guidatore}

\chapter{Implementazione del sistema}
\section{Architettura}
\section{Schema di connessione dei componenti}
\section{Componenti hardware utilizzati}
\section{Software e librerie utilizzate}
\section{Scelte progettuali}
\subsection{Filesystem}
\subsection{Sincronizzazione}
\section{Comunicazioni e trasmissione dati}
\subsection{OBD - esp32}
\subsection{esp32 - server remoto}

\chapter{Risultati ottenuti}
\section{Analisi dei risultati}
\section{Comportamenti non implementati}
\chapter{Conclusioni}
\section{Problemi aperti}
\section{Sviluppi futuri}


\hrule
%\bibliographystyle{plain}
\printbibliography[title={Bibliografia}]
%\addcontentsline{toc}{chapter}{Bibliografia}

\end{document}
